\section{Síntesis y ampliación}

\begin{importantbox}{Resumen de las ideas centrales de esta sesión}
\begin{enumerate}
  \item \textbf{Arquitectura}: MoE es una solución de ingeniería a corto plazo, Dense es la dirección a largo plazo; los FLOPs siempre son lo más barato.
  \item \textbf{Multimodalidad}: Native Multimodal se ve favorable a largo plazo, la separación entre LLM y VLM es un estado temporal. En la entrada no hay controversia sobre la unificación; en la salida todavía hay discrepancias.
  \item \textbf{Datos}: los datos sintéticos enfrentan riesgos de migración entre generaciones y de Model Collapse; Mid-training es el estado intermedio de entrenamiento hacia la modelización de un World Model. El autosupervisado Vision-Centric podría resolver el cuello de botella de datos multimodales.
  \item \textbf{RL}: la naturaleza de control en lazo cerrado hace que la definición del problema sea más importante que el algoritmo; la falta de una Scaling Law, la generalización insuficiente y la complejidad del environment son los tres grandes desafíos. El RL avanza hacia la especialización.
  \item \textbf{Agent}: la baja calidad de las decisiones y la gestión de Memory son el cuello de botella central; Claude Code es el referente de 2025; llevarlo directamente al mundo real para obtener la señal de reward podría ser la vía de ruptura.
  \item \textbf{Aprendizaje continuo}: el enfoque Context-based (Prompt Database + Retrieval) es más seguro y práctico que el Parameter-based; el contexto largo sigue siendo, en esencia, una cuestión de FLOPs.
\end{enumerate}
\end{importantbox}

\subsection{Lecturas adicionales}
\begin{itemize}
  \item Informe técnico de DeepSeek-R1: un nuevo paradigma de posentrenamiento con RL
  \item Entrada de blog de OpenAI ``From Compute-Bound to Data-Bound''
  \item Trabajos relacionados con Visual Jigsaw / Spatial Intelligence Benchmark
  \item Diseño de producto de Claude Code y arquitectura Agentic
  \item Bibliografía clásica sobre Slow Weight / Fast Weight (Hinton et al.)
\end{itemize}

\end{document}
